% Generated draft: requires mathematical and source review.
\documentclass[11pt]{article}
\usepackage[margin=25mm]{geometry}
\usepackage{fontspec}
\setmainfont{Latin Modern Roman}
\usepackage{amsmath,amssymb,mathtools,unicode-math}
\setmathfont{Latin Modern Math}
\usepackage{enumitem,needspace,xurl,hyperref}
\setlength{\parindent}{0pt}
\setlength{\parskip}{5pt}
\setlength{\emergencystretch}{3em}
\newcommand{\N}{\mathbb{N}}
\newcommand{\Z}{\mathbb{Z}}
\newcommand{\Q}{\mathbb{Q}}
\newcommand{\R}{\mathbb{R}}
\newcommand{\C}{\mathbb{C}}
\newcommand{\F}{\mathbb{F}}
\newcommand{\A}{\mathbb{A}}
\newcommand{\PP}{\mathbb{P}}
\newcommand{\E}{\mathbb{E}}
\newcommand{\eps}{\varepsilon}
\newcommand{\dd}{\,\mathrm{d}}
\begin{document}
\section*{UnsolvedMath ID: 30002709}
\textbf{Source problem code:} catalogue entry
\subsection{Definitions and mathematical statement}
Let \(L\subseteq F\) be fields with \([F:L]<\infty\), and let \(v:F\to\Gamma\cup\{\infty\}\) be a valuation, where \(\Gamma\) is an ordered abelian group. The valuation on \(L\) is the restriction of \(v\). For \(K\in\{L,F\}\), define the valuation ring and its maximal ideal by
\[
O_K=\{x\in K:v(x)\geq 0\},\qquad
\mathfrak m_K=\{x\in K:v(x)>0\}.
\]
For finitely many elements \(a_1,\ldots,a_r\in O_F\), put
\[
A=O_L[a_1,\ldots,a_r],\qquad \mathfrak p=A\cap\mathfrak m_F.
\]
Then \(\mathfrak p\) is a prime ideal of \(A\), and its localization, viewed inside \(F\), is
\[
A_{\mathfrak p}=\left\{\frac{b}{s}:b\in A,\ s\in A\setminus\mathfrak p\right\}.
\]
The supplied question asks whether, for every such finite valued-field extension \((F\mid L,v)\), there exist an integer \(r\geq 1\) and elements \(a_1,\ldots,a_r\in O_F\) such that
\[
O_F=\bigl(O_L[a_1,\ldots,a_r]\bigr)_{O_L[a_1,\ldots,a_r]\cap\mathfrak m_F}.
\]
Equivalently, is \(O_F\) always essentially of finite type over \(O_L\) in this specified sense? The supplied statement imposes no additional restriction on characteristic, valuation rank, separability, henselianity, or normality.

The supplied dated curation labels the problem partially solved and reports a known necessary condition, a theorem characterizing finite generation of the integral closure, and sufficiency for henselian or normal extensions. It does not state those conditions, identify the precise integral closure in that theorem, or specify the hypotheses and meaning of those sufficiency claims. These contextual claims therefore cannot be formulated precisely from this record and are not additional hypotheses of the displayed question.
\subsection{Short English statement}
For every finite extension of valued fields, can the larger valuation ring be obtained by adjoining finitely many of its elements to the smaller valuation ring and then localizing at the contraction of its maximal ideal? The supplied partially solved label and special-case claims are dated curation requiring independent review.
\subsection{Sources}
\begin{itemize}
\item Local UnsolvedMath registry: supplied complete record with permanent ID \(30002709\), title \emph{Essential Finite Generation of Valuation Rings}, source problem code \texttt{catalogue entry}, in the Oberwolfach Reports Open Problems collection. The proposed year is recorded as 2014; no individual proposer is supplied.
\item Supplied source citation: \emph{Valuation Theory and Its Applications} (2015), Oberwolfach Reports, pp.~2757--2823, DOI \texttt{10.4171/owr/2014/49}. \url{https://doi.org/10.4171/owr/2014/49}. The supplied literature list also identifies this as \emph{Oberwolfach Report 49/2014}: \url{https://publications.mfo.de/bitstream/handle/mfo/3439/OWR_2014_49.pdf?isAllowed=y&sequence=1}.
\item Supplied literature reference: Kuhlmann--Novacoski, \emph{Henselian elements}. \url{https://arxiv.org/abs/1311.6155}.
\item Literature assessment date: \texttt{2026-08-21}. The partially solved label and the assessment concerning necessary conditions, integral-closure finite generation, and sufficiency for henselian or normal extensions are supplied dated curation derived from the Oberwolfach material, not independently verified here. The registry also supplies the statement-status label \texttt{corrected\_verified} and reports that its clean statement was checked against the official report; that verification claim is supplied metadata, not an independent verification by this draft.
\end{itemize}
\end{document}
